\paragraph{反向二项核的加权正交闭式：对角自伴随权的唯一性、范数恒等式与谱投影条目化}
\begin{remark}[文献定位与陈述范围]\label{rem:xi-binomial-matrix-literature-scope}
关于反向二项核（Pascal 反序矩阵）在 \(\QQ(\sqrt5)\) 上的谱结构与对角化框架，参见 \cite{StanicaPeeleBinomialMatrix2000,LiverancePitsenberger1996DiagonalizationBinomialMatrix}。
本段仅在既有对角化框架之上补充三类与可审计实现直接对齐的闭式恒等式：对角自伴随权的唯一性、黄金本征向量族在二项权内积下的正交范数闭式、以及由此推出的谱投影与条目级重构公式。
\end{remark}
\par
固定整数 \(q\ge 2\)。
记反向二项核 \(B_q\in \Mat_{q+1}(\ZZ)\) 为
\[
(B_q)_{k\ell}:=\binom{q-\ell}{k}\qquad(0\le k,\ell\le q),
\]
并令二项权矩阵
\[
W_q:=\mathrm{diag}\Bigl(\binom{q}{0}^{-1},\binom{q}{1}^{-1},\dots,\binom{q}{q}^{-1}\Bigr),
\qquad
\langle u,v\rangle_{W_q}:=u^{\mathsf T}W_qv.
\]

\begin{conclusion}[二项权自伴随化的唯一性：对角权仅差整体比例]\label{con:xi-bq-unique-diagonal-selfadjoint}
设 \(D=\mathrm{diag}(d_0,\dots,d_q)\succ 0\) 为对角矩阵并满足
\[
B_q^{\mathsf T}D=DB_q.
\]
则 \(D\)（在正比例意义下）唯一确定，并且必为
\[
\boxed{\;
D\ \propto\ W_q
\;=\;
\mathrm{diag}\Bigl(\binom{q}{0}^{-1},\binom{q}{1}^{-1},\dots,\binom{q}{q}^{-1}\Bigr).\;}
\]
\end{conclusion}
\begin{proof}
将 \(B_q^{\mathsf T}D=DB_q\) 取 \((k,0)\) 元，得到
\[
(B_q^{\mathsf T}D)_{k0}=(B_q)_{0k}\,d_0=\binom{q-k}{0}\,d_0=d_0,
\qquad
(DB_q)_{k0}=d_k(B_q)_{k0}=d_k\binom{q}{k}.
\]
故对一切 \(0\le k\le q\) 有 \(d_k=d_0/\binom{q}{k}\)，从而 \(D\) 由 \(d_0\) 唯一确定并与 \(W_q\) 成正比例。证毕。
\end{proof}

\begin{conclusion}[黄金谱的正交范数闭式：线性型本征向量的 \texorpdfstring{$W_q$}{Wq}-规范化]\label{con:xi-bq-golden-spectrum-orthogonality-norm}
令 \(\varphi=\frac{1+\sqrt5}{2}\)。
对每个 \(j\in\{0,1,\dots,q\}\)，定义齐次多项式
\[
P_j(x,y):=(\varphi x+y)^{q-j}(x-\varphi y)^{j}
\,=\,\sum_{k=0}^{q}c^{(j)}_k\,x^{q-k}y^{k},
\qquad
c^{(j)}:=(c^{(j)}_0,\dots,c^{(j)}_q)^{\mathsf T}\in\RR^{q+1}.
\]
则 \(c^{(j)}\) 为 \(B_q\) 的本征向量且本征值为
\[
\boxed{\;\mu_j=(-1)^j\varphi^{\,q-2j}\;}
\qquad(0\le j\le q).
\]
并且在 \(\langle\cdot,\cdot\rangle_{W_q}\) 下满足严格正交与范数闭式
\[
\langle c^{(j)},c^{(j')}\rangle_{W_q}=0\quad(j\neq j'),
\qquad
\boxed{\;\langle c^{(j)},c^{(j)}\rangle_{W_q}=\frac{(\varphi+2)^{q}}{\binom{q}{j}}.\;}
\]
因此归一化向量
\[
u^{(j)}:=\sqrt{\frac{\binom{q}{j}}{(\varphi+2)^q}}\;c^{(j)}
\]
构成一组 \(W_q\)-正交规范基：\(u^{(j)\mathsf T}W_qu^{(j')}=\delta_{jj'}\)。
\end{conclusion}
\begin{proof}[证明要点]
\par\noindent\textbf{(i) 代入算子与本征关系}
对任意系数向量 \(c=(c_0,\dots,c_q)^{\mathsf T}\)，令 \(F_c(x,y):=\sum_{k=0}^{q}c_kx^{q-k}y^k\)。
由二项式展开可验证 \(B_q\) 对应代入作用 \((x,y)\mapsto(x+y,x)\)：
\[
F_{B_qc}(x,y)=F_c(x+y,x).
\]
对线性型 \(L_+(x,y):=\varphi x+y\)、\(L_-(x,y):=x-\varphi y\) 有
\[
L_+(x+y,x)=\varphi(x+y)+x=\varphi^2x+\varphi y=\varphi\,L_+(x,y),
\]
\[
L_-(x+y,x)=(x+y)-\varphi x=y-\varphi^{-1}x=-\varphi^{-1}L_-(x,y),
\]
其中使用 \(\varphi^2=\varphi+1\)。
故
\[
P_j(x+y,x)=\varphi^{q-j}(-\varphi^{-1})^{j}P_j(x,y)=(-1)^j\varphi^{q-2j}P_j(x,y),
\]
从而 \(B_qc^{(j)}=\mu_jc^{(j)}\)。

\par\noindent\textbf{(ii) 正交性}
由命题 \ref{prop:pom-bq-weighted-selfadjoint-binomial}，\(B_q\) 在 \(\langle\cdot,\cdot\rangle_{W_q}\) 下自伴随；
又由命题 \ref{prop:pom-bq-is-symq-and-spectrum}，本征值 \(\{\mu_j\}_{j=0}^{q}\) 两两不同。
因此不同本征值对应的本征向量两两正交，得到第一式。

\par\noindent\textbf{(iii) 范数闭式}
定义齐次次数 \(q\) 的 Fischer--apolar 配对
\[
\langle F,G\rangle_{\mathrm{ap}}
:=\frac{1}{q!}\,F(\partial_x,\partial_y)\,G(x,y)\big|_{(0,0)}.
\]
对基 \(x^{q-k}y^k\) 的直接计算给出
\(
\langle F_c,F_{c'}\rangle_{\mathrm{ap}}=c^{\mathsf T}W_qc'
\)，
故只需计算 \(\langle P_j,P_j\rangle_{\mathrm{ap}}\)。
令方向导数算子
\[
D_+:=\varphi\,\partial_x+\partial_y,\qquad D_-:=\partial_x-\varphi\,\partial_y.
\]
则 \(D_+(\varphi x+y)=\varphi^2+1=\varphi+2\)、\(D_+(x-\varphi y)=0\)，以及 \(D_-(\varphi x+y)=0\)、\(D_-(x-\varphi y)=\varphi+2\)。
因此 \(D_+\) 仅作用于第一因子而 \(D_-\) 仅作用于第二因子，得到
\[
D_+^{\,q-j}D_-^{\,j}P_j
=(q-j)!\,j!\,(\varphi+2)^{q}.
\]
再由定义
\(
\langle P_j,P_j\rangle_{\mathrm{ap}}
=\frac{1}{q!}D_+^{\,q-j}D_-^{\,j}P_j
\)
得到
\(
\langle c^{(j)},c^{(j)}\rangle_{W_q}
=\frac{(q-j)!\,j!}{q!}(\varphi+2)^q=(\varphi+2)^q/\binom{q}{j}
\)。
证毕。
\end{proof}

\begin{conclusion}[谱投影与逆本征矩阵闭式：无需求逆即可条目级重构 \texorpdfstring{$B_q^{\,m}$}{Bq^m}]\label{con:xi-bq-spectral-projector-entrywise}
令
\[
C_q:=\bigl[c^{(0)},c^{(1)},\dots,c^{(q)}\bigr]\in\RR^{(q+1)\times(q+1)},
\qquad
\Lambda_q:=\mathrm{diag}(\mu_0,\dots,\mu_q),
\]
并记 \(W_q\) 如上。
则 \(C_q\) 可逆且其逆矩阵满足闭式
\[
\boxed{\;
C_q^{-1}=(\varphi+2)^{-q}\,W_q^{-1}\,C_q^{\mathsf T}\,W_q\;}
\]
并由此对任意整数 \(m\ge 0\) 有谱投影展开
\[
\boxed{\;
B_q^{\,m}
=(\varphi+2)^{-q}\,C_q\,\Lambda_q^{\,m}\,W_q^{-1}\,C_q^{\mathsf T}\,W_q\; } .
\]
特别地，每个条目 \((k,\ell)\) 具有完全显式有限和表示
\[
\boxed{\;
\bigl(B_q^{\,m}\bigr)_{k\ell}
=\frac{1}{(\varphi+2)^{q}\binom{q}{\ell}}
\sum_{j=0}^{q}\binom{q}{j}\,\mu_j^{m}\,c^{(j)}_{k}\,c^{(j)}_{\ell}\; } .
\]
\end{conclusion}
\begin{proof}
由结论 \ref{con:xi-bq-golden-spectrum-orthogonality-norm} 的正交范数闭式，矩阵式可写为
\[
C_q^{\mathsf T}W_qC_q=(\varphi+2)^q\,W_q.
\]
左乘 \(W_q^{-1}\) 得 \(W_q^{-1}C_q^{\mathsf T}W_qC_q=(\varphi+2)^q I\)，从而
\(
C_q^{-1}=(\varphi+2)^{-q}W_q^{-1}C_q^{\mathsf T}W_q
\)。
由对角化 \(B_q=C_q\Lambda_qC_q^{-1}\) 得到所示 \(B_q^{\,m}\) 的谱投影展开；条目公式由右端逐项展开并注意 \(W_q\) 的对角性读出。证毕。
\end{proof}

\begin{conclusion}[完备性核恒等式：一条“黄金--二项”的离散分辨公式]\label{con:xi-bq-resolution-identity}
沿结论 \ref{con:xi-bq-spectral-projector-entrywise} 的记号，取 \(m=0\) 得对任意 \(0\le k,\ell\le q\) 成立恒等式
\[
\boxed{\;
\sum_{j=0}^{q}\binom{q}{j}\,c^{(j)}_{k}\,c^{(j)}_{\ell}
=(\varphi+2)^{q}\binom{q}{\ell}\,\delta_{k\ell}\; } .
\]
\end{conclusion}
\begin{proof}
将结论 \ref{con:xi-bq-spectral-projector-entrywise} 的条目公式取 \(m=0\)，并用 \((B_q^0)_{k\ell}=\delta_{k\ell}\) 即得。证毕。
\end{proof}

\endinput
